\documentclass[11pt,a4paper]{article}
\usepackage[T1]{fontenc}
\usepackage{lmodern}
\usepackage[margin=1in]{geometry}
\usepackage{amsmath,amssymb,amsthm}
\usepackage[hidelinks]{hyperref}
\newtheorem{theorem}{Theorem}
\newtheorem{lemma}[theorem]{Lemma}
\newcommand{\R}{\mathbb{R}}
\newcommand{\Reach}{\mathcal{R}}
\newcommand{\dd}{\,\mathrm{d}}
\title{Failure of finite-jet determination for small-time local controllability}
\author{Kaiwen Jin\\\small\href{mailto:yui.ui.siki@gmail.com}{yui.ui.siki@gmail.com}}
\date{October 15, 2026}

\begin{document}
\maketitle

\begin{abstract}
Small-time local controllability asks whether a system can reach a
neighborhood of an equilibrium in arbitrarily short time.
We give a negative answer to the finite-jet determination question for
real-analytic systems with bounded measurable controls.
Starting from a previously constructed five-dimensional polynomial system
with two controls, we show that adding a positive polynomial term of
arbitrarily high order destroys this property.
The original system is fixed independently of the prescribed jet order.
The proof compares a logarithmic bound for closed trajectories with the
positive integral cost of an auxiliary coordinate's excursion.
\end{abstract}

\section{The finite-jet question and the example}

Consider real-analytic vector fields $X_0,\ldots,X_m$ near $0\in\R^n$,
with $X_0(0)=0$, and the control-affine system
\[
 \dot x=X_0(x)+\sum_{j=1}^m u_j(t)X_j(x),\qquad x(0)=0,
 \qquad u(t)\in[-1,1]^m.
\]
Controls are measurable and trajectories are absolutely continuous
solutions of the equations almost everywhere.
Write $\Reach_{\le t}(0)$ for the endpoints reachable in a duration at most
$t$. The system is small-time locally controllable at zero if
$\Reach_{\le t}(0)$ contains a Euclidean neighborhood of zero for every
sufficiently small $t>0$.

The finite-jet question asks whether, for each such controllable system,
there is an integer $N$ such that every real-analytic tuple with the same
Taylor coefficients through order $N$ at zero is again small-time locally
controllable with the same control set.
This is Conjecture~1.1 in the v7 preprint of Jafarpour~\cite{Jafarpour}.
His Theorem~7.1 proves finite-jet stability under a polynomial lower bound
on reachable neighborhoods.
Jin~\cite{Jin} constructed a polynomial system that is small-time locally
controllable without such a lower bound.
We show that the same system has no finite jet guaranteeing controllability.
The system and the logarithmic estimate used below are from~\cite{Jin};
the additional argument controls the positive cost of a high-order
perturbation for every admissible control.

Let $e_j$ denote the $j$th standard coordinate vector in $\R^5$.
The original system is
\begin{equation}\label{eq:original}
\begin{aligned}
 \dot x_1&=u, & \dot x_2&=x_1, & \dot x_3&=x_2,\\
 \dot x_4&=x_1^2x_3^2-\left(\frac19+x_5^2\right)x_2^4,
 &\dot x_5&=v,
\end{aligned}
\end{equation}
with measurable $|u|,|v|\le1$.
Denote its drift by $X_0$ and its controlled fields by $X_1=e_1$ and
$X_2=e_5$. All three fields are polynomial and $X_0(0)=0$.

\begin{theorem}\label{thm:main}
The fixed system~\eqref{eq:original} is small-time locally controllable at
zero. For every integer $N\ge0$, there is an integer $k\ge1$ with $2k>N$
such that the polynomial tuple
\begin{equation}\label{eq:perturbation}
 Y_0=X_0+x_5^{2k}e_4,\qquad Y_1=X_1,\qquad Y_2=X_2
\end{equation}
is not small-time locally controllable at zero with the same control set
$[-1,1]^2$, although
\[
 D^\alpha Y_j(0)=D^\alpha X_j(0)
 \qquad (0\le j\le2,\ |\alpha|\le N).
\]
Here $\alpha$ is a multi-index and $|\alpha|$ its total order.
This five-dimensional example disproves the universal finite-jet
implication for real-analytic systems in dimensions at least three.
\end{theorem}

The five-dimensional example suffices for this universal question.
No separate conclusion in dimensions three or four is asserted.
The perturbation in~\eqref{eq:perturbation} adds a nonnegative term to the
fourth equation. Its integral is positive whenever the fifth coordinate
leaves zero, and this contribution will dominate every possible negative fourth
displacement in sufficiently small time.

\section{Positive perturbations of arbitrarily high order}

Theorem~1 of~\cite{Jin} proves the controllability of
\eqref{eq:original}. Its proof constructs trajectories that return the first three and fifth
coordinates to zero while producing either sign of fourth displacement,
and then steers the other four coordinates,
reaching a full neighborhood in a duration at most $5t/8$ for each $t>0$.
We also use the following estimate, stated here with its full hypotheses.

\begin{lemma}[{\cite[Lemma~2]{Jin}}]\label{lem:log}
Let $T>0$ and let $f\in C^2([0,T])$ have Lipschitz second derivative,
$\|f'''\|_{L^\infty}\le1$, and
\[
 f(0)=f'(0)=f''(0)=f(T)=f'(T)=f''(T)=0.
\]
Set
\[
 A=\int_0^T f'^4\dd s,\qquad
 E=\int_0^T\left(ff''+\frac13f'^2\right)^2\dd s.
\]
Then
\begin{equation}\label{eq:critical}
 E=\int_0^T f^2f''^2\dd s-\frac19A.
\end{equation}
If $A>0$, then $A\le T^9/16$ and
\begin{equation}\label{eq:log}
 E\ge\frac{A}{36\log^2(T^9/(4A))}.
\end{equation}
\end{lemma}

The identity~\eqref{eq:critical} follows by integration by parts.
The quantitative lower bound~\eqref{eq:log} is the logarithmic estimate
proved in~\cite{Jin}. We retain the original normalization of $f,A,E$.

\begin{proof}[Proof of Theorem~\ref{thm:main}]
The original system's controllability is the cited Theorem~1 of~\cite{Jin}.
Fix an integer $k\ge1$ and consider the perturbed tuple
\eqref{eq:perturbation}. Its first, second, third and fifth equations are
unchanged, while
\[
 \dot y_4=y_1^2y_3^2-\left(\frac19+y_5^2\right)y_2^4+y_5^{2k}.
\]
Successive integration gives a unique trajectory on every finite interval
for each bounded measurable control. In particular, no existence or
finite-time escape restriction enters the argument.
We use the usual integral calculus for absolutely continuous functions
\cite[Appendix~3.A, Example~3.59 and Theorems~3.60, 3.62]{Hunter}.

Take any control of duration $T>0$ whose trajectory ends on the fourth
axis, and put $f=y_3$ and $g=y_5$.
The unchanged equations give $f'=y_2$, $f''=y_1$ and $f'''=u$ almost
everywhere. Thus $f$ satisfies every hypothesis of Lemma~\ref{lem:log},
including the six zero endpoint jets.
Also $g(0)=g(T)=0$ and $g$ is $1$-Lipschitz, since $g'=v$ almost everywhere.
Define
\[
 a=\max_{0\le s\le T}|g(s)|\le T,\qquad
 P_k=\int_0^T g(s)^{2k}\dd s.
\]
The fourth displacement is exactly
\begin{equation}\label{eq:displacement}
 y_4(T)=E-\int_0^T g^2f'^4\dd s+P_k,
\end{equation}
by~\eqref{eq:critical}.

If $a=0$, this displacement is $E\ge0$.
If $A=0$, continuity gives $f'=0$, hence $f=0$, $E=0$, and
$y_4(T)=P_k\ge0$.
If $E\ge a^2A$, then~\eqref{eq:displacement} also gives $y_4(T)\ge0$.
It remains to consider $a>0$, $A>0$ and $E<a^2A$.
Since $A\le T^9/16$, the logarithm in~\eqref{eq:log} is positive.
That estimate now gives
\[
 \frac{1}{36\log^2(T^9/(4A))}\le\frac EA<a^2,
\]
and therefore
\begin{equation}\label{eq:small-A}
 A<\frac{T^9}{4}\exp\left(-\frac1{6a}\right).
\end{equation}

The excursion in $g$ also forces a positive lower bound for $P_k$.
Choose $s_*\in[0,T]$ with $|g(s_*)|=a$.
The Lipschitz bound and $g(0)=0$ imply $s_*\ge a$ and
\[
 |g(s)|\ge a-(s_*-s)\qquad (s_*-a\le s\le s_*).
\]
This estimate holds whether $g(s_*)$ is positive or negative. Hence
\begin{equation}\label{eq:cost}
 P_k\ge\int_{s_*-a}^{s_*}(a-s_*+s)^{2k}\dd s
       =\frac{a^{2k+1}}{2k+1}.
\end{equation}
Combining~\eqref{eq:displacement}, \eqref{eq:small-A} and
\eqref{eq:cost}, we obtain
\begin{equation}\label{eq:comparison}
 y_4(T)>\frac{a^{2k+1}}{2k+1}
          -\frac{T^9a^2}{4}\exp\left(-\frac1{6a}\right).
\end{equation}
The positive power series for the exponential gives
\[
 \exp\left(-\frac1{6a}\right)
 \le (2k)!6^{2k}a^{2k}.
\]
Writing $C_k=(2k)!6^{2k}$ and using $a\le T$, the comparison becomes
\[
 y_4(T)>a^{2k+1}\left(\frac1{2k+1}-\frac{C_kT^{10}}4\right).
\]
Set
\begin{equation}\label{eq:tau}
 \tau_k=\min\left\{1,
       \left(\frac{2}{(2k+1)C_k}\right)^{1/10}\right\}>0.
\end{equation}
For $0<T\le\tau_k$, the expression in parentheses is at least
$1/(2(2k+1))>0$.
Together with the previous cases, this proves that every fourth-axis
endpoint in every such duration has nonnegative fourth coordinate.

For each $0<t\le\tau_k$, the union of endpoints over all durations
$0\le T\le t$ therefore excludes every point $-r e_4$ with $r>0$.
Every neighborhood of zero contains such a point, so the perturbed system
is not small-time locally controllable.

Finally, given any finite $N\ge0$, choose $k\ge1$ with $2k>N$.
The only changed field component is $x_5^{2k}$ in the drift.
Every partial derivative of this monomial of total order at most $N$
vanishes at zero: a derivative in another coordinate is identically zero,
and derivatives only in $x_5$ retain a positive power of $x_5$.
The controlled fields are unchanged and $Y_0(0)=0$.
The tuple $X$ has been fixed throughout, completing the asserted
finite-jet counterexample.
\end{proof}

\paragraph{AI-use statement.}
OpenAI Codex was used for mathematical development, source checks, drafting
and internal manuscript review.

\begin{thebibliography}{9}
\bibitem{Jafarpour}
S.~Jafarpour,
\emph{On small-time local controllability},
SIAM Journal on Control and Optimization \textbf{58} (2020), 425--446.
\href{https://doi.org/10.1137/16M1068797}{doi:10.1137/16M1068797}.
The conjecture and theorem locators refer to the
\href{https://arxiv.org/abs/1604.02432v7}{v7 preprint} (2019).

\bibitem{Jin}
K.~Jin,
\emph{Small-time controllability without polynomial reachable neighborhoods},
research note, October 12, 2026.
\href{https://github.com/yuiyuiui/AIM/blob/fff30eeb072395341d9944e4c75d5267b770f3b8/research/solutions/417-polynomial-reachable-growth-counterexample/paper/main.pdf}{Versioned public note},
Theorem~1, Lemma~2 and Section~3.

\bibitem{Hunter}
J.~K.~Hunter,
\emph{Notes on Partial Differential Equations}, Appendix~3.A,
\href{https://www.math.ucdavis.edu/~hunter/pdes/ch3A.pdf}{Absolutely continuous functions},
Example~3.59 and Theorems~3.60, 3.62, p.~79.
\end{thebibliography}
\end{document}
